% ============================================================
% Round 07: DAPO
% ============================================================

\subsection{Round 07：DAPO}

\subsubsection{本轮问题}

\begin{conceptbox}
\textbf{学习目标：} 理解 DAPO 是如何在 GRPO 基础上改进大规模 reasoning RL 训练稳定性的。重点掌握四个组件：Clip-Higher、Dynamic Sampling、Token-Level Policy Gradient Loss、Overlong Reward Shaping。\\
\textbf{主线位置：}
\[
\text{GRPO}
\to
\text{DAPO}
\to
\text{Dr.GRPO / GSPO / REINFORCE++}
\]
\end{conceptbox}

\begin{conceptbox}
\textbf{Highlight：DAPO}\\
\textbf{论文：} \emph{DAPO: An Open-Source LLM Reinforcement Learning System at Scale}\\
\textbf{arXiv：} \href{https://arxiv.org/abs/2503.14476}{2503.14476}，submitted 2025-03-18，last revised 2025-05-20\\
\textbf{Project：} \href{https://dapo-sia.github.io/}{dapo-sia.github.io}\\
\textbf{Code：} \href{https://github.com/BytedTsinghua-SIA/DAPO}{BytedTsinghua-SIA/DAPO}\\
\textbf{Data / Model：} \href{https://huggingface.co/datasets/BytedTsinghua-SIA/DAPO-Math-17k}{DAPO-Math-17k}，
\href{https://huggingface.co/BytedTsinghua-SIA/DAPO-Qwen-32B}{DAPO-Qwen-32B}\\
\textbf{关键词：} GRPO improvement，long-CoT RLVR，Clip-Higher，Dynamic Sampling，Token-Level Loss，Overlong Reward Shaping\\
\textbf{核心定位：} 在 GRPO 基础上补齐大规模 reasoning RL 的工程 recipe，让训练更稳定、更可复现。
\end{conceptbox}

\subsubsection{DAPO 是什么}

DAPO = Decoupled Clip and Dynamic sAmpling Policy Optimization。

它不是从零发明一个新 RL 框架，而是在 GRPO 的基础上，针对 long-CoT reasoning RL 的几个工程痛点做系统修补。

标准 GRPO 的核心是：

\[
\text{同题采样一组回答}
\to
\text{组内 reward normalization}
\to
\text{PPO-style ratio + clip 更新 policy}
\]

DAPO 仍然保留这条主线，但加入四个关键技巧：

\begin{enumerate}[leftmargin=1.5em]
  \item Clip-Higher；
  \item Dynamic Sampling；
  \item Token-Level Policy Gradient Loss；
  \item Overlong Reward Shaping。
\end{enumerate}

\begin{summarybox}
\textbf{一句话：} DAPO = 更工程化、更稳定的 GRPO recipe，目标是让 long-CoT RLVR 真正跑起来、跑得稳、跑得有效。
\end{summarybox}

\subsubsection{DAPO 为什么出现}

论文里的出发点很直接：naive GRPO 在复现大规模 reasoning RL 时会遇到：

\begin{itemize}[leftmargin=1.5em]
  \item entropy collapse：模型分布越来越尖，探索变少；
  \item reward noise：尤其是超长被截断样本带来的噪声；
  \item training instability：训练曲线不稳，性能上不去；
  \item sample inefficiency：很多 group 全对或全错，几乎没有梯度。
\end{itemize}

所以 DAPO 的四个组件分别对应这些问题：

\begin{center}
{\small
\begin{tabular}{p{3.6cm}p{4.2cm}p{5cm}}
\hline
\textbf{组件} & \textbf{修的问题} & \textbf{核心动作} \\
\hline
Clip-Higher
&
entropy collapse，探索不足
&
提高 ratio 的上 clip，让好样本概率有更多上升空间
\\
Dynamic Sampling
&
全对 / 全错 group 没梯度
&
只保留同题内既有正确又有错误的 group
\\
Token-Level Loss
&
sample-level loss 对长回答权重不合理
&
按所有有效 token 汇总 loss，而不是每条回答先平均
\\
Overlong Reward Shaping
&
截断样本 reward 噪声
&
过滤或软惩罚超长回答，避免粗暴惩罚污染训练
\\
\hline
\end{tabular}
}
\end{center}

\subsubsection{DAPO 的目标函数}

可以把 DAPO 看成 GRPO objective 的一个改进版：

\[
\mathcal{J}_\text{DAPO}(\theta)
=
\mathbb{E}
\left[
\frac{1}{
\sum_{i=1}^{G}|o_i|
}
\sum_{i=1}^{G}
\sum_{t=1}^{|o_i|}
\min
\left(
\rho_{i,t}(\theta)\hat{A}_{i,t},\;
\operatorname{clip}
\left(
\rho_{i,t}(\theta),
1-\epsilon_\text{low},
1+\epsilon_\text{high}
\right)
\hat{A}_{i,t}
\right)
\right]
\]

其中：

\[
\rho_{i,t}(\theta)
=
\frac{
\pi_\theta(o_{i,t}\mid q,o_{i,<t})
}{
\pi_{\theta_\text{old}}(o_{i,t}\mid q,o_{i,<t})
}
\]

\[
\hat{A}_{i,t}
\approx
\hat{A}_i
=
\frac{
r_i-\operatorname{mean}(\{r_j\}_{j=1}^{G})
}{
\operatorname{std}(\{r_j\}_{j=1}^{G})+\epsilon
}
\]

注意这里写成 $\hat{A}_{i,t}$，但在常见 outcome reward 设置里，同一个回答 $o_i$ 的所有 token 通常共享同一个回答级 advantage $\hat{A}_i$。

\begin{warnbox}
\textbf{和标准 GRPO 的两个明显差别：}
\begin{itemize}[leftmargin=1.5em]
  \item clip 从对称的 $[1-\epsilon,1+\epsilon]$ 变成非对称的 $[1-\epsilon_\text{low},1+\epsilon_\text{high}]$；
  \item loss normalization 从每条回答内部先除以 $|o_i|$，变成对整个 group 的有效 token 总数归一化。
\end{itemize}
\end{warnbox}

\subsubsection{追问：DAPO 为什么公式里没有 KL}

前面 PPO / GRPO 里我们经常写 KL penalty：

\[
-
\beta
\mathbb{D}_\text{KL}
(\pi_\theta \| \pi_\text{ref})
\]

它的作用是防止 policy 偏离 reference model 太远。

但 DAPO 论文在 long-CoT reasoning RL 里做了一个重要取舍：\key{去掉显式 KL penalty}。直觉是，reasoning RL 的目标不是让模型一直贴近初始 instruct model，而是允许模型发展出更长、更复杂的推理分布。

\[
\text{long-CoT RL}
\quad
\Rightarrow
\quad
\pi_\theta \text{ 可能需要明显偏离 } \pi_\text{ref}
\]

如果 KL 约束太强，可能会压制这种推理能力的展开。

\begin{warnbox}
\textbf{注意：} 这不是说所有 GRPO / RLVR 训练都必须去 KL。是否使用 KL、KL 系数设多大，是一个工程超参数。DAPO 的经验是在它的 long-CoT 数学 RL 设置里，去掉显式 KL 有利于训练。
\end{warnbox}

\subsubsection{组件一：Clip-Higher}

标准 PPO / GRPO 的 clip 是对称的：

\[
\operatorname{clip}(\rho,1-\epsilon,1+\epsilon)
\]

DAPO 改成：

\[
\operatorname{clip}
(\rho,1-\epsilon_\text{low},1+\epsilon_\text{high})
\]

并且通常设置：

\[
\epsilon_\text{high} > \epsilon_\text{low}
\]

例如论文中使用过：

\[
\epsilon_\text{low}=0.20,
\qquad
\epsilon_\text{high}=0.28
\]

为什么要这样做？

如果一个 token 原本概率很低，但是它出现在高 reward 的探索性回答里，标准对称 clip 可能很快限制它的概率上升。这样模型会更容易强化已有高概率 token，而不是探索低概率但有价值的新路径。

\begin{intubox}
\textbf{直觉：} Clip-Higher 给“提高好样本概率”更多空间，尤其是低概率探索 token，从而缓解 entropy collapse。
\end{intubox}

它不是放开所有更新。下界仍然保持较小：

\[
1-\epsilon_\text{low}
\]

这样可以避免过度降低某些 token 概率，导致采样空间快速坍缩。

\subsubsection{组件二：Dynamic Sampling}

GRPO 的学习信号来自组内差异。假设同一个 prompt 采样 $G=4$ 个回答。

如果 reward 是：

\[
[1,1,1,1]
\]

或者：

\[
[0,0,0,0]
\]

那么组内 advantage 基本都是 $0$，这组样本几乎不会提供有效梯度。

DAPO 的 Dynamic Sampling 直接过滤这种 group。它保留满足下面条件的 group：

\[
0
<
\left|
\{o_i \mid \operatorname{is\_correct}(o_i)\}
\right|
<
G
\]

也就是：

\begin{itemize}[leftmargin=1.5em]
  \item 至少有一个正确回答；
  \item 至少有一个错误回答。
\end{itemize}

\begin{summarybox}
\textbf{Dynamic Sampling = 只训练有区分度的 group。}\\
全对没有可比差异，全错也没有正样本方向；混合 group 才能告诉模型“同一道题里，哪些路径比哪些路径更值得强化”。
\end{summarybox}

这会多花一些采样成本，但减少无效训练 batch，整体可能更快收敛。

\subsubsection{组件三：Token-Level Policy Gradient Loss}

原始 GRPO 常见写法是 sample-level loss：

\[
\mathcal{L}_\text{GRPO}
=
-
\frac{1}{G}
\sum_{i=1}^{G}
\frac{1}{|o_i|}
\sum_{t=1}^{|o_i|}
l_{i,t}
\]

这个写法先对每条回答内部取平均，再对回答取平均。结果是：

\[
\text{每条回答的总权重差不多一样}
\]

不管回答有多长。

在 long-CoT 场景里，这会带来问题：

\begin{itemize}[leftmargin=1.5em]
  \item 高质量长推理里的每个 token 权重被稀释，学得不够；
  \item 低质量超长回答里的重复、乱码、绕圈 token 也被稀释，惩罚不够；
  \item 模型可能学到不健康的长度增长。
\end{itemize}

DAPO 改成 token-level loss：

\[
\mathcal{L}_\text{DAPO}
=
-
\frac{1}{
\sum_{i=1}^{G}|o_i|
}
\sum_{i=1}^{G}
\sum_{t=1}^{|o_i|}
l_{i,t}
\]

也就是所有有效 token 放在一起平均。

\begin{intubox}
\textbf{直觉：} sample-level loss 是“每篇作文权重一样”；token-level loss 是“每个有效 token 权重一样”。在长 CoT 里，后者更能处理长回答中的具体模式。
\end{intubox}

\subsubsection{组件四：Overlong Reward Shaping}

RL 训练时通常会设置最大生成长度：

\[
T_\text{max}
\]

如果回答超过最大长度，就会被截断。

问题是：被截断不一定代表推理是坏的。有时模型可能正在进行合理推理，但因为太长被截断。如果直接给所有截断样本一个强负 reward，会引入 reward noise。

DAPO 讨论了两类做法：

\begin{itemize}[leftmargin=1.5em]
  \item \textbf{Overlong Filtering：} 对被截断样本 mask loss，不让它们污染训练；
  \item \textbf{Soft Overlong Punishment：} 不一刀切，而是在超过安全长度后逐渐增加惩罚。
\end{itemize}

可以抽象写成：

\[
R_\text{final}
=
R_\text{accuracy}
+
R_\text{length}
\]

其中 $R_\text{length}$ 在回答过长时变成负数，并随长度逐渐变大。

\begin{warnbox}
\textbf{关键点：} DAPO 不是简单惩罚所有长回答。它要避免“长但合理的推理”被粗暴打成坏样本，同时也要防止模型无限延长、重复和乱码。
\end{warnbox}

\subsubsection{DAPO 和 GRPO 的关系}

\begin{center}
{\small
\begin{tabular}{p{3.2cm}p{4.7cm}p{4.9cm}}
\hline
 & \textbf{GRPO} & \textbf{DAPO} \\
\hline
Advantage
&
组内 reward normalization
&
基本沿用组内 advantage
\\
Clip
&
对称 clip：$1-\epsilon,1+\epsilon$
&
非对称 clip：$1-\epsilon_\text{low},1+\epsilon_\text{high}$
\\
采样
&
普通 group sampling
&
dynamic sampling，过滤全对 / 全错 group
\\
Loss 聚合
&
常见 sample-level mean
&
token-level mean
\\
超长样本
&
容易带来截断 reward 噪声
&
filtering 或 soft overlong punishment
\\
目标
&
去 critic 的 PPO-style RL
&
让 long-CoT RLVR 训练更稳、更有效
\\
\hline
\end{tabular}
}
\end{center}

\subsubsection{一个完整 DAPO iteration}

可以把 DAPO 的一轮训练想成：

\begin{enumerate}[leftmargin=1.5em]
  \item 从题库采样一批 prompts；
  \item 令 $\pi_{\theta_\text{old}} \leftarrow \pi_\theta$；
  \item 对每个 prompt 用 old policy 采样 $G$ 个回答；
  \item 用 rule-based verifier 给每个回答打 reward；
  \item 过滤全对 / 全错 group，保留有区分度的 group；
  \item 对保留的 group 计算组内 advantage；
  \item 处理超长样本：过滤或加 soft penalty；
  \item 用当前 policy 重新计算这些回答 tokens 的 logprob；
  \item 计算 ratio，并使用 asymmetric clip；
  \item 用 token-level loss 更新 policy。
\end{enumerate}

\begin{summarybox}
\textbf{最小结论：} DAPO 不是一个完全替代 GRPO 的新范式，而是一个更强的 GRPO 工程配方。\\
它的四个关键词是：
\[
\text{Clip-Higher}
\quad
\text{Dynamic Sampling}
\quad
\text{Token-Level Loss}
\quad
\text{Overlong Shaping}
\]
\end{summarybox}

\subsubsection{本轮检查题}

\begin{enumerate}[leftmargin=1.5em]
  \item DAPO 为什么要使用 Clip-Higher，而不是标准对称 clip？
  \item Dynamic Sampling 为什么要过滤全对和全错的 group？
  \item sample-level loss 和 token-level loss 的区别是什么？
  \item 为什么被截断的超长样本会带来 reward noise？
  \item DAPO 和 GRPO 的关系是“替代”还是“增强”？为什么？
\end{enumerate}
